\chapter{Vectors and Lists}
\label{chap:vectors-lists}

\chapterguide
  {You should be able to create variables, call functions, and use logical/relational operators.}
  {construct vectors and lists; access elements using position, logical, negative, and name-based indexing; perform vector arithmetic; recognize vector recycling in the supplied example; sort vectors; name list elements; and add, remove, or update list entries.}

\section{R vectors}

Lab 3a begins with vector construction. The examples emphasize that a vector holds a sequence of values and can be generated either with the colon operator, \texttt{seq()}, or \texttt{c()}.

\subsection{Constructing vectors}

\begin{lstlisting}[style=dsr]
# Integer-like sequence
v <- 5:13
print(v)

# Sequence beginning at a decimal value
v <- 6.6:12.6
print(v)

# End value does not fall exactly on the sequence
v <- 3.8:11.4
print(v)

# Explicit step size
print(seq(5, 9, by = 0.4))
\end{lstlisting}

The handout also demonstrates type coercion by mixing text, a number, and a logical value in one \texttt{c()} call:

\begin{lstlisting}[style=dsr]
s <- c('apple', 'red', 5, TRUE)
print(s)
\end{lstlisting}

The supplied note states that the numeric and logical values are converted to characters in this mixed vector.

\begin{keyidea}
In the handout's mixed-vector example, all elements end up represented as character values. This is one reason the course later introduces lists: lists can keep different element types together without forcing them into a single common representation.
\end{keyidea}

\subsection{Accessing vector elements}

Lab 3a gives four indexing patterns using a days-of-week vector:

\begin{lstlisting}[style=dsr]
t <- c("Sun", "Mon", "Tue", "Wed", "Thurs", "Fri", "Sat")

# Position indexing
u <- t[c(2, 3, 6)]
print(u)

# Logical indexing
v <- t[c(TRUE, FALSE, FALSE, FALSE, FALSE, TRUE, FALSE)]
print(v)

# Negative indexing
x <- t[c(-2, -5)]
print(x)

# 0/1 indexing example from the handout
y <- t[c(0, 0, 0, 0, 0, 0, 1)]
print(y)
\end{lstlisting}

These examples show different ways to express the same general operation: choose which positions should be returned.

\section{Vector arithmetic}

The handout creates two vectors of equal length and applies arithmetic directly:

\begin{lstlisting}[style=dsr]
v1 <- c(3, 8, 4, 5, 0, 11)
v2 <- c(4, 11, 0, 8, 1, 2)

print(v1 + v2)
print(v1 - v2)
print(v1 * v2)
print(v1 / v2)
\end{lstlisting}

The operations are applied element by element. The division example includes division by a zero element in \texttt{v2}; observe R's output rather than assuming every element will be a finite ordinary number.

\subsection{Vector element recycling}

Lab 3a next shortens the second vector:

\begin{lstlisting}[style=dsr]
v1 <- c(3, 8, 4, 5, 0, 11)
v2 <- c(4, 11)

# Handout note: v2 behaves as c(4, 11, 4, 11, 4, 11)
print(v1 + v2)
print(v1 - v2)
\end{lstlisting}

The exercise explicitly explains the repeated form of \texttt{v2}. This is the vector-recycling behavior the lab wants you to notice.

\section{Sorting vectors}

\begin{lstlisting}[style=dsr]
v <- c(3, 8, 4, 5, 0, 11, -9, 304)
print(sort(v))
print(sort(v, decreasing = TRUE))

v <- c("Red", "Blue", "yellow", "violet")
print(sort(v))
print(sort(v, decreasing = TRUE))
\end{lstlisting}

The same \texttt{sort()} function is used for numeric and character vectors, with \texttt{decreasing = TRUE} reversing the order requested.

\section{R lists}

A list is introduced using a mixture of strings, numbers, a vector, and a logical value:

\begin{lstlisting}[style=dsr]
list_data <- list(
  "Red",
  "Green",
  c(21, 32, 11),
  TRUE,
  51.23,
  119.1
)
print(list_data)
\end{lstlisting}

This contrasts with the mixed vector: the list is explicitly intended to hold heterogeneous components.

\subsection{Naming list elements}

\begin{lstlisting}[style=dsr]
list_data <- list(
  c("Jan", "Feb", "Mar"),
  list("green", 12.3)
)

names(list_data) <- c("1st Quarter", "A Inner list")
print(list_data)
\end{lstlisting}

Names make later access more readable than relying only on positions.

\subsection{Accessing list elements}

The handout uses numeric position and name-based access:

\begin{lstlisting}[style=dsr]
list_data <- list(
  c("Jan", "Feb", "Mar"),
  list("green", 12.3)
)
names(list_data) <- c("1st_Quarter", "A_Inner_list")

print(list_data[1])
print(list_data[2])
print(list_data$A_Inner_list)
print(which(list_data$`1st_Quarter` == "Feb"))
\end{lstlisting}

The \texttt{which()} line combines list-name access with a relational test to find the position of \texttt{"Feb"}.

\subsection{Manipulating a list}

\begin{lstlisting}[style=dsr]
list_data <- list(
  c("Jan", "Feb", "Mar"),
  list("green", 12.3)
)
names(list_data) <- c("1st_Quarter", "A_Inner_list")

# Add a third element
list_data[3] <- "New_element"
print(list_data[3])

# Remove it
list_data[3] <- NULL
print(list_data[3])

# Update the second element
list_data[2] <- "updated_element"
print(list_data[2])
\end{lstlisting}

The exercise intentionally uses the same indexing form for adding and updating, and assigns \texttt{NULL} to remove an element.

\begin{quickcheck}
\begin{enumerate}
  \item What happens to the numeric and logical values in the handout's mixed character vector?
  \item How does negative vector indexing differ from positive position indexing?
  \item What repeated form of \texttt{v2} is stated in the vector-recycling example?
  \item Why is a list suitable for the Lab 3a example containing strings, numbers, a vector, and a logical value?
  \item How does \texttt{list\_data\$A\_Inner\_list} differ in intent from accessing an unnamed position?
\end{enumerate}
\end{quickcheck}

\begin{chapterrecap}
\begin{itemize}
  \item Vectors are built with \texttt{:}, \texttt{seq()}, and \texttt{c()} in the supplied exercises.
  \item Vector indexing can be position-based, logical, negative, or expressed through the 0/1 example in the handout.
  \item Arithmetic on equal-length vectors is element-wise; the handout also demonstrates recycling of a shorter vector.
  \item \texttt{sort()} orders numeric and character vectors and can reverse the direction.
  \item Lists are used for heterogeneous components and can be named, accessed, extended, updated, and shortened.
\end{itemize}
\end{chapterrecap}
